\begin{abstract}
Scientific data discovery requires search that thinks before it acts. An effective system must profile each corpus to understand what metadata exist, decide which search strategies apply, execute domain-specific reasoning such as the arithmetic needed to recognize that 200~kPa, 200000~Pa, and 0.2~MPa describe identical pressure, and refine its approach when results are insufficient. This observe-decide-act-evaluate cycle is what distinguishes agentic search from conventional retrieval. Few current systems perform these operations in combination.

We present CLIO Search, an agentic retrieval system for scientific data. CLIO implements a four-stage pipeline: Discover (profile each corpus in 14--22\,ms), Reason (select retrieval branches based on corpus characteristics), Transform (execute dimensional conversion and formula normalization as deterministic database operators alongside BM25 and dense vector search), and Execute (federate across storage backends with iterative refinement via an LLM rewriter or deterministic fallback achieving 96\% of LLM-assisted quality).

On cross-unit queries, CLIO achieves 0.99~P@5 where BM25 scores 0.40 and NC-Retriever achieves 0.16; on same-unit controls all methods score comparably (0.30--0.45). Precision improves by 87\% on 288 real NOAA documents. At distributed scale on 2.5M arXiv abstracts across 4~NVIDIA GH200 nodes, cross-unit queries return identical results at sub-50\,ms latency with 0.94 weak scaling efficiency. All code and benchmarks are open-source.
\end{abstract}

%% ===================================================================
